\chapter{El plan que nombra a otro municipio}

\label{cap:pdua}
\section{El plan provincial ya lo había dicho}

Antes de los instrumentos municipales que este capítulo analiza, la provincia se dio un plan
propio. Y ese plan dice, sobre el crecimiento del área metropolitana, lo mismo que este libro
documenta expediente por expediente.

\begin{ficha}{Plan de Desarrollo Estratégico de la Provincia de Salta 2030}
Aprobado por el \textbf{Decreto 2478/12} en los términos del artículo 77 de la Constitución
Provincial, «que determina el carácter definitorio del Plan».

\textbf{Origen:} iniciativa privada de la Fundación Salta (2009), con metodología del IAE
Business School de la Universidad Austral y subsidio del Banco Interamericano de Desarrollo.
Treinta y siete mesas de concertación, quinientas treinta y siete personas.

En el diagnóstico del área de ordenamiento territorial:

«El crecimiento provincial ha fluido espontáneamente, \textbf{sin planificación}, lo que generó
desequilibrios que no resultan fácilmente abordables. Uno de los problemas que se ha tornado
visible es el \textbf{dinámico crecimiento del área metropolitana} y la despoblación de las
zonas rurales, además de situaciones varias que deben enfrentarse para \textbf{paliar los
efectos de la falta de previsión}.»

Y entre las amenazas del diagnóstico social: «\textbf{Falta de accesibilidad del Estado a
tierras aptas para la urbanización}».
\end{ficha}

\textbf{Primera: el Estado provincial declaró en 2012 que el crecimiento metropolitano ocurrió
sin planificación.} Es anterior a todos los loteos que este libro documenta y a los dos
instrumentos municipales que este capítulo analiza. \textbf{Lo que el libro reconstruye acto
por acto, el plan provincial lo enunció como diagnóstico antes de que ocurriera la mayor parte
de esos actos.}

\textbf{Segunda: la falta de tierra pública apta figura como amenaza declarada.} Es,
exactamente, el segundo circuito que el capítulo \ref{cap:tierrafiscal} documenta: un Estado que adjudica por
sorteo lo poco que tiene mientras el circuito privado opera sin ese límite.

\textbf{Tercera: el plan tiene carácter definitorio por mandato constitucional.} No es un
documento indicativo: el decreto que lo aprueba invoca el artículo 77 de la Constitución
provincial. Este libro no encontró ningún acto del corredor norte que lo invoque.

\textbf{Y cuarta: participaron los actores que reaparecen después.} En el Consejo de Control y
Seguimiento está el \textbf{COPAIPA} --- el mismo cuerpo técnico que la jueza reconoce en
2026 por su labor en el Plan de Manejo ---; en las mesas participaron la \textbf{Intendencia
de La Caldera}, la \textbf{Autoridad Metropolitana de Transporte} y el \textbf{Instituto
Provincial de los Pueblos Indígenas}. \textbf{Los cuatro reaparecen en este libro catorce años
después, en expedientes que ninguno de ellos conectaba entonces.}

\begin{cautela}
Este libro leyó el plan completo ---ciento setenta páginas sobre trece sectores--- pero no pudo establecer qué se ejecutó de él. El corredor norte no figura entre sus proyectos: «La Caldera» aparece dos veces en todo el documento, según detalla el capítulo \ref{cap:trabajo}.

Lo que registra es que \textbf{el diagnóstico existe, es oficial, es anterior, y coincide}.
\end{cautela}

\section{La norma marco}

La Caldera aprobó su Plan de Desarrollo Urbano Ambiental por la Ordenanza Nº 643 del 25 de noviembre de 2015. El plan consta de dos piezas: un Informe Final y un Proyecto de Ordenanza elevado al Concejo Deliberante. \textbf{Este libro leyó ese proyecto, no el texto sancionado}: si la Ordenanza 643 lo aprobó sin modificaciones ---y por lo tanto si el artículo 6º que los capítulos \ref{cap:fincas} y \ref{cap:tierra} citan es el del texto vigente--- lo dice el libro de actas del Concejo. \pendiente{texto de la Ordenanza Nº 643 tal como se sancionó}

\begin{ficha}{Ficha documental --- PDUA del Municipio de La Caldera}
\begin{itemize}[leftmargin=*,nosep]
\item \textbf{Comitente:} Consejo Federal de Inversiones. \textbf{Ejecutores:} Fundación del Colegio de Arquitectos y Subsecretaría de Planeamiento Urbano de la Provincia de Salta.
\item \textbf{Equipo:} coordinador urbano Arq. Durbal Basualdo; coordinador de infraestructura Ing. Hugo Calvo; Arq. Silvia Ale; Arq. Mariana Zoricich; Geól. Raúl Chocobar.
\item \textbf{Agradecimientos:} Arq. Adriana Krumpholz (Subsecretaría provincial); Arq. Ignacio Pancetti (Fundación del Colegio de Arquitectos); Municipio de La Caldera; Ing. Virginia del Val.
\item \textbf{Metodología declarada:} relevamiento in situ, planos de diagnóstico, concertación mediante audiencia pública y presentación de resultados en una segunda audiencia.
\item \textbf{Ámbito:} ejido municipal de La Caldera y paraje de La Calderilla. Mantiene en vigencia las ordenanzas 040/82 y 243/97.
\end{itemize}
\end{ficha}

Hay continuidad técnica real entre el plan y su reglamentación posterior: la Arq. Zoricich integra el equipo de 2015 y es coautora y coordinadora de edición del Código de Ordenamiento Territorial; la Arq. Krumpholz aparece en ambos documentos como enlace de la Subsecretaría provincial.

\subsection*{Lo que el Estado provincial dijo del plan al presentarlo}

El parte de prensa del Ministerio de Gobierno del 10 de diciembre de 2015 confirma la ordenanza y agrega cuatro datos que el propio informe no consigna.

\begin{ficha}{Presentación oficial, 10 de diciembre de 2015}
Informó la Arq.\ \textbf{Adriana Krumpholz}, entonces secretaria de Planeamiento Urbano del Ministerio de Gobierno. Consignó que el trabajo en la localidad \textbf{comenzó en 2013}; que La Caldera es el \textbf{cuarto municipio} de la provincia en obtener aprobación de un plan de este tipo; y que los puntos centrales son la recuperación de áreas verdes, la definición de sectores de preservación, la \textbf{delimitación de terrenos para vivienda social sólo para habitantes del pueblo} y la delimitación de áreas de crecimiento urbano \textbf{con proyección de quince a veinte años}. Consignó además que \textbf{se proyecta la construcción de una costanera} que descongestione la Ruta 9, con espacios de esparcimiento y recreación.
\end{ficha}

\begin{cautela}
Tres lecturas.

\textbf{Que sea el cuarto municipio} explica el error de plantilla que se describe más abajo, sin excusarlo. El plan no es una pieza única redactada para La Caldera: es la cuarta aplicación de un dispositivo serial. Un texto que se reutiliza es un texto donde el nombre del municipio anterior sobrevive si nadie lo relee, y eso es exactamente lo que ocurrió.

\textbf{Que el trabajo comenzara en 2013} sitúa la elaboración en el mismo año en que la Provincia declaraba, por el DNU 1595/13 convalidado como Ley 7781, la emergencia hídrica en todo el territorio con orden de readecuar y ordenar todas las cuencas. El plan se elaboró bajo emergencia hídrica declarada, y dejó el estudio de cota máxima edificable como proyecto pendiente.

\textbf{Y la costanera es la pieza que más pesa.} El plan de 2015 proyecta una costanera sobre el río. En febrero de 2019, el escrito del amparo colectivo enumera entre las zonas críticas las actividades en la playa del río sobre la \textbf{avenida Costanera sur}, cuyos desbordes alcanzaron a los loteos del sur. No se afirma aquí que la costanera proyectada y la avenida Costanera del amparo sean la misma obra, ni que exista relación causal entre el proyecto y el daño: eso requiere el plano del plan y un peritaje, y ninguno de los dos se consultó. Deja registrado que el instrumento que zonificó el crecimiento proyectó una obra vial sobre la ribera del mismo río cuya línea de ribera a lo largo del pueblo la Provincia había fijado en febrero de 2014, y cuyos afluentes ---los arroyos El Durazno y Guaranguay, con línea fijada en diciembre de ese año--- inundarían los loteos cuatro años después.
\end{cautela}

\subsection*{El programa tiene nombre, y lo financiaba el BID}

El parte de 2015 decía que La Caldera era el cuarto municipio en obtener un plan de este tipo, sin nombrar a los anteriores. Un decreto de 2014 nombra uno de los programas que producían esos planes, y el municipio vecino al que se aplicó.

\begin{ficha}{Decreto Nº 397 del 11 de febrero de 2014 --- B.O.\ Nº 19.249}
Expediente 0110272-199.922/2013-0. Aprueba la contratación del consultor individual \textbf{coordinador del Proyecto ``Plan Urbano Ambiental (PUA) para el Municipio de Vaqueros''}, en el marco del \textbf{``Programa de Mejora de la Gestión Municipal'', Préstamo BID 1855/OC-AR}.

Consultor: \textbf{arquitecto Luis Fernando Eladio Kuncar}. La Unidad Ejecutora Provincial es la Subsecretaría de Financiamiento; la Unidad Ejecutora Central, del Ministerio del Interior de la Nación.

El dictamen jurídico deja constancia de que, por haberse cumplido las pautas del Banco Interamericano de Desarrollo, \textbf{la contratación queda excluida del régimen de la Ley 6.838} de contrataciones de la provincia, conforme su artículo 97 inciso a.

Firman: Urtubey, Parodi, Simón Padrós.
\end{ficha}

Tres consecuencias.

\textbf{Primera: el instrumento de Vaqueros existe desde 2014 y tiene autor.} Este libro venía consignando el PIDUA-VAQ como pendiente sin fecha ni origen. Su coordinación se contrató en febrero de 2014, con financiamiento internacional, casi dos años antes de que se aprobara el plan de La Caldera.

\textbf{Segunda: los planes no son iniciativas municipales.} Son productos de programas financiados desde afuera del municipio: en Vaqueros, un préstamo del BID con unidad ejecutora en el Ministerio del Interior de la Nación; en La Caldera, el Consejo Federal de Inversiones. El municipio recibe el plan; no lo encarga ni lo paga. Eso vuelve verosímil la serialidad que el error de plantilla, descripto más abajo, deja ver, sin atribuirla a desidia local, aunque la plantilla de La Caldera no proviene de este préstamo.

\textbf{Tercera, y es la incómoda: la contratación queda fuera del régimen provincial de compras.} No por excepción discrecional sino por el artículo 97 inciso a de la propia Ley 6.838, que excluye lo financiado por organismos internacionales cuando se siguen sus pautas. Es legal y es habitual. Pero significa que el equipo que redactó el plan urbano de Vaqueros ---por un circuito del mismo tipo que el que produjo el de La Caldera, cuyo informe nombra a Rivadavia Banda Sur--- fue seleccionado por un procedimiento que no rinde cuentas ante el sistema de contrataciones de la provincia sino ante el manual de un banco.

\begin{cautela}
Y el contrato de consultoría del plan de Vaqueros también se publicó. La \textbf{Resolución Nº 409 D del 8 de junio de 2015} aprueba el contrato celebrado el 6 de mayo entre la Subsecretaría de Financiamiento y la firma \textbf{AGEMET S.R.L.}, en el marco de la Solicitud de Propuesta 05/14, para la \textbf{``Elaboración de la Planificación Integral de Desarrollo Urbano Ambiental del Municipio de Vaqueros''}, dentro del mismo Programa de Mejora de la Gestión Municipal y del mismo préstamo BID 1855/OC-AR. Monto: \textbf{\$1.830.000}.

El expediente documenta un procedimiento competitivo real: invitaciones publicadas en diarios de circulación nacional y local, apertura de propuestas técnicas en agosto de 2014 con las económicas reservadas y cerradas, evaluación técnica con no objeción de la unidad ejecutora central, y recién después apertura de precios. Compitieron \textbf{Serman Asociados} (75 puntos técnicos), \textbf{Oficina Urbana S.A.} (84) y \textbf{AGEMET S.R.L.} (86), y se descartó la propuesta de Marcelo Arzelan \& Asociados por exceder el presupuesto fijo del pliego.

Conviene decirlo con precisión, porque matiza lo que este capítulo observó antes. La contratación quedó fuera del régimen de la Ley 6.838, pero \textbf{no fue discrecional}: siguió el manual del banco, con publicidad, comisión evaluadora, sobres cerrados y puntaje técnico previo al precio. Lo que este libro señaló --- que rinde cuentas ante el manual de un banco y no ante el sistema provincial de contrataciones --- sigue siendo cierto y sigue importando, pero no equivale a ausencia de competencia.

Lo que no cambia es la escala del programa: un millón ochocientos treinta mil pesos por el plan de un municipio, financiados por un préstamo internacional, con la firma seleccionada por la Subsecretaría de Financiamiento provincial, la no objeción de la unidad ejecutora central y el municipio como destinatario. El plan de La Caldera, aprobado meses después, salió de un circuito análogo ---comitente federal, ejecución provincial y el municipio como destinatario---, aunque su financiamiento fue del Consejo Federal de Inversiones y no del préstamo del BID.
\end{cautela}


\section{Lo que el artículo 6º está enumerando}

El proyecto de ordenanza enumera, en su artículo 6º, las localidades y parajes que integran el municipio; el capítulo \ref{cap:fincas} lee esa lista como mapa de la tenencia histórica.

Leída en 2015, esa lista parece una descripción geográfica. El capítulo \ref{cap:fincas} muestra que es otra
cosa. \textbf{El Gallinato, San Alejo, Los Porongos y Potrero de Castilla son cuatro de los
trece partidos policiales en que la provincia dividió el departamento en 1909}; La Calderilla
y Yacones también. La Caldera identifica su propio territorio, en una norma de planeamiento
del siglo XXI, con las unidades administrativas de hace un siglo.

No hay nada impropio en eso: los topónimos duran, y una ordenanza que enumera parajes debe
usar los nombres que la gente usa. Pero conviene tenerlo presente al leer el resto de este
capítulo, porque el plan que va a fijar zonificaciones y usos del suelo está trabajando sobre
un mapa que hereda sin discutir. \textbf{Cuando el artículo 6º dice «Finca», está diciendo
una unidad de dominio}, y este libro no pudo establecer, para ninguna de las seis, quién es
hoy su titular.

\section{Rivadavia Banda Sur}

La introducción del Informe Final dice, textualmente, que el Consejo Federal de Inversiones es el ente que encomienda el Plan de Desarrollo Urbano Ambiental \textbf{para el municipio de Rivadavia Banda Sur}. No es un desliz aislado: el apartado final, que lista los proyectos complementarios sugeridos, los presenta como contemplados en el sistema de usos de suelo propuesto \textbf{para Rivadavia Banda Sur}.

Rivadavia Banda Sur es un municipio del Chaco salteño, a más de cuatrocientos kilómetros de La Caldera, en el departamento que el libro anterior documentó como frente de desmonte y territorio del caso Lhaka Honhat. No tiene yungas, ni cornisa, ni dique.

\begin{cautela}
Conviene medirlo con precisión.

Lo que \textit{no} prueba: no prueba que el plan sea una copia. El cuerpo del informe es específico de La Caldera --- describe el río, la cornisa, el dique, La Calderilla, el casco histórico, la ruta 9 como única vía de acceso --- y los planos son del territorio caldereño. El trabajo de campo existió.

Lo que \textit{sí} establece: el documento fue producido sobre una plantilla de otro municipio, y esa plantilla sobrevivió sin corregir en la apertura y en el cierre, es decir en los dos lugares que cualquier lector abre primero. Pasó por el equipo redactor, por la Fundación del Colegio de Arquitectos, por la Subsecretaría de Planeamiento Urbano provincial, por el CFI como comitente, por el Ejecutivo municipal, por dos audiencias públicas y por el Concejo Deliberante que lo aprobó como ordenanza. Ninguna de esas instancias lo detectó, o ninguna lo consideró suficiente para devolver el documento.

La conclusión no es sobre la honestidad de nadie. Es sobre la naturaleza del control: un instrumento de planificación que llega a ordenanza municipal conservando el nombre de otro municipio en su primer párrafo no fue leído íntegramente por ninguna de las instancias que debían aprobarlo.
\end{cautela}

\section{El diagnóstico que el plan hace de sí mismo}

Descontado el error, el diagnóstico del informe es lúcido y es la voz oficial describiendo el proceso que este libro documenta.

Sobre las áreas vacantes, consigna que en ellas se identifican sectores de uso agropecuario y, en algunos casos, \textbf{emprendimientos privados de loteos en ejecución sin planificación urbana}. Sobre las causas, sostiene que la falta de planificación y la gran demanda de terrenos generan asentamientos privados y edificaciones espontáneas en lugares poco convenientes. Sobre la estructura social, describe dos polos de crecimiento: los pobladores locales y los nuevos habitantes que inmigran en general al sector de La Calderilla. Sobre el riesgo, señala que el principal problema del ejido es la accesibilidad, con una única vía estrecha, y que las áreas de expansión llegan hasta los límites naturales \textbf{a definir}, según la cota máxima de edificación y las líneas de ribera.

El informe reproduce además el diagnóstico del Plan de Ejecución Metropolitano de 2012 sobre los municipios del conurbano salteño: absorbidos por la expansión del área central, que les modifica su función; y advierte que, si no se interviene previamente, se transformarán en periferia degradada \textbf{o en base de asentamiento de sectores cerrados}, rompiendo la tradición cultural y local de esas localidades. Consigna que en el período intercensal previo los mayores crecimientos del área metropolitana correspondieron a La Caldera, con 37,3 por ciento, y Cerrillos, con 35,2, mientras el resto creció entre 13 y 15.

\begin{cautela}
El municipio incorporó a su norma, en 2015, la descripción anticipada de lo que le iba a ocurrir. La advertencia sobre la conversión en base de asentamiento de sectores cerrados está escrita en los fundamentos del instrumento que debía impedirlo. Diez años después, el Concejo Deliberante pide la clausura preventiva del ingreso a un loteo cerrado.
\end{cautela}

\section{La cota que nunca se determinó}

La fórmula ``límites naturales a definir'' se repite en el informe cada vez que se menciona el borde de la expansión. Y el primer punto de la lista de proyectos complementarios sugeridos es: estudios necesarios para la determinación de la cota máxima edificable.

\begin{cautela}
El plan que fija dónde puede crecer La Caldera declara, en su propio texto, que el dato técnico que define ese límite todavía no existe y debe encargarse por separado. La expansión residencial se zonificó primero; el estudio que determina hasta qué cota es seguro construir quedó como proyecto pendiente.

Tres años después de aprobado el plan, los aluviones de diciembre de 2018 y enero de 2019 afectaron a los loteos del sur. Nada permite afirmar que exista relación causal --- eso lo dirá un peritaje hidrológico, no un ensayo --- pero deja registrado que el estudio de cota máxima edificable figuraba como pendiente en el instrumento de planificación vigente al momento del daño.
\end{cautela}


\section{Lo que el proyecto de ordenanza fija y lo que deja abierto}

El proyecto clasifica el territorio en zona urbana, sub-urbana y rural; subdivide el área urbana en consolidada, a consolidar y de expansión; define quince puntos de coordenadas para el polígono urbano y veinticinco para el sub-urbano; y establece once distritos: R1 y M1, R2 y AS, ASE, R3 y R4, más AE, AT, APN y RN. En R3 y R4 la superficie mínima de parcela es de 600 m² con quince metros de frente, retiro de tres metros y ocupación máxima del cincuenta por ciento.

Dos ausencias merecen registro.

\textbf{Las cesiones.} Los fundamentos afirman que se han establecido para cada zona las condiciones de fraccionamiento del suelo. El articulado no las contiene: no hay en todo el proyecto un porcentaje de cesión obligatoria de espacio verde ni institucional, ni fórmula para calcularlos.

\textbf{La discrecionalidad.} La fórmula ``quedará a consideración del municipio'' aparece de manera recurrente: para asentamientos residenciales o turísticos en zona sub-urbana y rural, para construcciones en área especial, para cambios de uso en área de protección natural, para edificaciones en área de turismo, para excepciones de altura y de retiro. No se establece qué debe evaluar el municipio, con qué dictamen técnico, en qué plazo ni con qué recurso para el afectado.

El Código de Ordenamiento Territorial posterior corrige esta segunda falla creando el Órgano Técnico de Aplicación y el procedimiento de evaluación de impacto. Pero entre 2015 y la sanción del Código, en 2016, rigió una zonificación precisa con una cláusula de excepción abierta; el trámite del agua del loteo El Durazno, iniciado en 2014, se había publicado en julio de 2015, meses antes de que el plan se aprobara; su plano se aprobaría en 2021, con el Código ya vigente.
